% ===== §14 =====
\section{The Polyphonic Conclusion}
\label{sec:polyphony}

\noindent
It remains to close, and the manner of the closing is not incidental to the argument but its last step. The paper has denied that intimate relation is an optimization problem, criticized the imposition of the frame, and recovered a rationality of response, attunement, sustaining, and non-action. A reader who has followed this far will feel the pull of a single question, and the whole of the close turns on how it is answered: if not utility, then \emph{what} should one optimize? Love, perhaps, or meaning, or generativity itself? The temptation to answer is strong, and it is strong for a reason worth naming, since the reader has been trained by the very frame the paper opposes to expect that a serious argument ends by supplying a new objective to put in the place of the one it rejected. An answer would round the paper off in the shape the frame has taught us to find satisfying, replacing the discarded objective with a worthier one and leaving the machinery of maximization intact but redirected. The paper refuses the answer, and the refusal is not an omission or a failure of nerve but the final content of its view.

\subsection*{The surrender concealed in naming a new objective}

The temptation must be met head-on, because to yield to it would undo everything the paper has argued, and would undo it quietly, wearing the look of a constructive conclusion. To say ``do not optimize utility; optimize love'' is to keep the entire apparatus the paper has dismantled, the objective, the scale, the maximum, the optimizer, and merely to change the argument of the function. But the four impossibilities were not impossibilities about \emph{utility} in particular. They were impossibilities about the optimizing form as such, brought to bear on intimate relation: that its goods resist representation, admit no exterior standpoint to fix them, share no common scale, and belong to no subject stable enough to own the objective. Not one of these failures is touched by a change of objective. Love, written as the argument of $J$, is as non-representable, as relational, as incommensurable, and as owned-by-no-stable-subject as any other relational good; to make it the thing maximized is to inflict upon it precisely the damage the criticism described, rendering the circulation a settleable quantity and drawing it off in the reckoning. The offer to optimize love rather than utility is therefore not a friendlier version of the optimizing frame but the same frame with a more sympathetic mask, and it does the same harm to a more beloved object; indeed it does the harm more surely, since the reverence in which love is held makes the reckoning feel like devotion and disarms the suspicion that would have attended the reckoning of mere utility. To accept the offer would be to lose the argument at the last step, having won it at every prior one. The paper therefore declines to name a master objective, not from an inability to think of one, but from having understood that any objective so named would reinstate the error the whole paper exists to expose.

\subsection*{The plural and the unsynthesizable}

What stands in the place where a master objective would go is not a blank but a positive thesis, and it is the thesis the series has reached before by other roads: that the goods and the rationalities of relation are plural, and that their plurality admits no synthesis into a single measure without loss. There is no one thing that relational rationality maximizes, and this is not because the one thing has yet to be found but because there is none to find. There are, rather, many goods standing to one another in relations no common scale represents, and many aspects of the reasoning that honors them, response and attunement and sustaining and restraint, that do not reduce to a single operation performed on a single objective. This is the counterpart, in the theory of rationality, of the polyphonic form the series has given its conclusions since the paper on sustainability, where it was argued that the truth of a matter may lie in the tension among several frameworks none of which subsumes the others, and that to force them into one synthesis is to lose exactly the truth that lived in their irreducibility \citep{paperix}. The refusal to name a master objective is that same refusal carried to the level of value: an insistence that the many goods and the many modes of their honoring are not to be collapsed into one, and that the collapse, however it satisfies the desire for a single answer, destroys what it unifies. The plurality is not a stage on the way to a unity not yet achieved. It is the finding.

\subsection*{The form as the argument}

That the paper enacts this refusal rather than merely asserting it is the last thing to say, and it is what binds the paper's form to its content in the manner the series requires of itself. A paper that argued against the single measure and then, in its conclusion, supplied one, would refute itself in its form whatever its sentences maintained; the shape of its ending would betray the burden of its argument, offering with the closing gesture exactly the master objective the body had shown to be impossible. This paper ends instead by declining the measure, and the declining is meant to be read as the argument's completion and not its evasion. The most honest thing a work against optimization can do at its close is to withhold the new optimum the reader has been trained to expect, to leave the plurality unsynthesized on the page, and to let the absence of a master objective be felt as the presence of the thesis. The reader who reaches the end and finds no single answer to ``what, then, should one optimize?'' has not been failed by a paper that lost its nerve before the constructive task. He has been given the answer, which is that the question, in this domain, has no answer of the kind it asks for, and that to have felt the want of one and let it stand unfilled is to have understood. The form withholds the summit because there is none, and the withholding is what the paper had to say.
